\SourceLeaf{290}%
\[
\begin{aligned}
\mathrm{d}\mathord{.}\frac{\operatorname{tang}\psi}{m}&=-\frac{l\mathrm{d}x}{m}\\
\mathrm{d}u&=\operatorname{tang}\psi\mathord{.}\mathrm{d}x
\end{aligned}
\]
we obtain
\[
\begin{aligned}
\operatorname{tang}\psi&=\mathfrak{A}(1+\mu x^3+\mu'x^4+\text{etc.})-\tfrac{1}{3}\lambda x^3-\tfrac{1}{4}\lambda'x^4-\text{etc.}\\
u&=\mathfrak{B}+\mathfrak{A}(x+\tfrac{1}{4}\mu x^4+\tfrac{1}{5}\mu'x^5+\text{etc.})-\tfrac{1}{12}\lambda x^4-\tfrac{1}{20}\lambda'x^5-\text{etc.}
\end{aligned}
\]
\indent The constants introduced by the integration, \(\mathfrak{A},\mathfrak{B}\), can be determined from the condition that \(u=0\) for the two values of \(x\) corresponding to the points \(F,G\). Let these values be \(x=-\tfrac{1}{2}(h-\delta)\) and \(x=+\tfrac{1}{2}(h+\delta)\), where \(\delta\) denotes the value of \(2x\) at the point midway between \(F\) and \(G\) and is generally a quantity of the same order as \(h\), or of a higher order if \(H\) lies very close to that midpoint. From this it is easy to derive the following expression, accurate through order \(h^3\) inclusive, for \(\mathfrak{A}\)
\[
\mathfrak{A}=\frac{\lambda((h+\delta)^4-(h-\delta)^4)}{192h}=\tfrac{1}{24}\lambda\delta(hh+\delta\delta)
\]
\indent Substitute this into the series for \(\operatorname{tang}\psi\). Then give the variable \(x\) the particular values \(-\tfrac{1}{2}(h-\delta),+\tfrac{1}{2}(h+\delta)\), and we obtain, again with third-order accuracy,
\[
\begin{aligned}
\operatorname{tang}\psi^0&=\tfrac{1}{24}\lambda h(hh-2h\delta+3\delta\delta)\\
\operatorname{tang}\psi'&=-\tfrac{1}{24}\lambda h(hh-2h\delta+3\delta\delta)
\end{aligned}
\]
\indent In the special case of the investigation to be developed later, the condition stated above is accompanied by a further circumstance. The standard parallel lies midway between the two parallels on which the points \(F,G\) lie. Consequently, the abbreviated expressions
\[
\begin{aligned}
\operatorname{tang}\psi^0&=\tfrac{1}{24}\lambda h^3\\
\operatorname{tang}\psi'&=-\tfrac{1}{24}\lambda h^3
\end{aligned}
\]
will already be accurate to third order, as is easy to show in the following way. Denote the latitude of \(F\) by \(Q+q\) and that of \(G\) by \(Q-q\). The spherical triangles (\(F,H\), pole) and (\(G,H\), pole) then give the equations
